% Fate's Edge SRD −−− Equipment (Section 06)

\section{Equipment}

\subsection{Weapons by Weight Class}

A weapon's weight determines its Harm, its dice modifier at different ranges, and its tempo. Harm is applied directly when an attack hits; the dice modifier affects the attack roll, not damage.

\begin{tabular}{@{}lccccl@{}}
⊤rule
\textbf{Weight} & \textbf{Harm} & \textbf{Close} & \textbf{Near} & \textbf{Far} & \textbf{Tempo} \\
∣rule
Light   & 1 & +2d & +1d & −−− & Fast \\
Medium  & 2 & +1d & +2d & −−− & Standard \\
Heavy   & 3 & –1d & +3d & +1d & Slow \\
⊥tomrule
\end{tabular}

⫕section{Notes}
\begin{itemize}
  ℹtem \textbf{Fast tempo:} You may move and attack in the same turn (Move + Action).
  ℹtem \textbf{Standard tempo:} You may move or attack in a turn, not both.
  ℹtem \textbf{Slow tempo:} Attacking requires a Set action (or a full turn action) before the attack; you cannot move and attack in the same turn unless a talent says otherwise.
  ℹtem \textbf{Heavy weapons at Close:} Suffer –1d due to limited maneuvering room.
  ℹtem A weapon cannot attack beyond its maximum listed range unless a talent or special ammunition extends it.
\end{itemize}

\subsection{Ranged Weapons and Ammunition}

Ranged weapons follow the same weight and tempo rules. Ammunition is tracked abstractly as part of Supply (see Section~\ref{sec:travel−framework}). A character may run low on ammunition when the GM spends SB on resource strain or during extended travel.

\subsection{Armor}

Armor does not reduce Harm directly. Instead, it converts incoming Harm into Fatigue. See Section~\ref{sec:core−mechanic} for the full conversion procedure.

\begin{tabular}{@{}lccc@{}}
⊤rule
\textbf{Armor Type} & \textbf{Conversion} & \textbf{Notes} \\
∣rule
Light   & 1 Harm $→$ 1 Fatigue & –1d to physical actions (optional) \\
Medium  & 2 Harm $→$ 1 Fatigue & No penalty \\
Heavy   & 3 Harm $→$ 2 Fatigue & –2d to physical actions; cannot sprint \\
⊥tomrule
\end{tabular}

⫕section{Armor Conversion Examples}
\begin{itemize}
  ℹtem Light armor vs Harm 1: take 1 Fatigue. No Harm.
  ℹtem Medium armor vs Harm 2: take 1 Fatigue. No Harm.
  ℹtem Heavy armor vs Harm 3: take 2 Fatigue. No Harm.
  ℹtem Minimum 1 Fatigue per hit, even if conversion would give 0.
\end{itemize}

\subsection{Shields (Optional)}

Shields are worn in the off−hand. Using a shield restricts weapon size to Light or Medium.

\begin{tabular}{@{}lcl@{}}
⊤rule
\textbf{Shield} & \textbf{Benefit} & \textbf{Tradeoff} \\
∣rule
Buckler    & +1d to Defend actions once per scene & Off−hand occupied; –1d to ranged attacks \\
Heater     & +1d to Defend actions; 1 Harm $→$ Fatigue per scene & Off−hand occupied; –1d to ranged attacks \\
Pavise     & Plant to create heavy cover (one direction) & Bulky; cannot move while planted \\
⊥tomrule
\end{tabular}

\subsection{Mundane Gear}

Most gear is abstracted. When you need a specific item (rope, lockpicks, torch), assume you have it unless the fiction says otherwise or the GM spends SB to declare a shortage.

The following items are tracked as part of \textbf{Supply} (see Section~\ref{sec:travel−framework}):

\begin{multicols}{2}
\begin{itemize}
  ℹtem Rations (1 day)
  ℹtem Torches (3)
  ℹtem Waterskin
  ℹtem Rope (10m)
  ℹtem Tinderbox
  ℹtem Bedroll
  ℹtem Crowbar
  ℹtem Hammer and spikes
  ℹtem Sack (large)
  ℹtem Grappling hook
  ℹtem Lantern and oil
  ℹtem Healer's kit (3 uses)
\end{itemize}
\end{multicols}

\subsection{Equipment Condition}

All gear can become \textbf{Neglected} or \textbf{Compromised}.

\begin{itemize}
  ℹtem \textbf{Neglected:} Weapons suffer –1 die on attack rolls; armor conversion worsens by one step (e.g., Light behaves as Unarmored). Neglected gear is usable but unreliable.
  ℹtem \textbf{Compromised:} Weapons suffer –1 die on the first attack each round; armor provides no conversion (all Harm applies directly). Compromised gear is barely functional.
\end{itemize}

⫕section{Repair}
\begin{itemize}
  ℹtem \textbf{Neglected $→$ Functional:} Requires a short rest (1 hour) and access to basic tools.
  ℹtem \textbf{Compromised $→$ Neglected:} Requires a scene of dedicated work (Craft test, DV 3) or a visit to a smith.
  ℹtem \textbf{Compromised $→$ Functional:} Requires significant downtime (a full day) and a successful Craft test (DV 4) or a professional smith.
  ℹtem Enchanted items require the same repairs; enchantments are inactive while the item is Compromised.
\end{itemize}

\subsection{Magical Items (Enchantments)}

Enchantments are permanent magical modifications treated as Minor or Standard Assets (see Section~\ref{sec:followers−assets}). Each enchantment costs XP and counts toward the item's enchantment limit.

⫕section{Enchantment Limits}
\begin{itemize}
  ℹtem Light item: maximum 1 enchantment.
  ℹtem Medium item: maximum 2 enchantments.
  ℹtem Heavy item: maximum 3 enchantments.
\end{itemize}

⫕section{Sample Enchantments}
\begin{itemize}
  ℹtem \textbf{Keen Edge (2 XP):} +1 die when attacking armored targets.
  ℹtem \textbf{Flaming Blade (4 XP):} +1 Harm against creatures vulnerable to fire; +1 Effect against cold−based foes.
  ℹtem \textbf{Soulfire Weapon (6 XP):} Ignores 1 point of armor or resistance.
  ℹtem \textbf{Shadowweave Armor (2 XP):} +1 die to Stealth rolls while moving.
  ℹtem \textbf{Runed Plate (4 XP):} Reduce magical backlash severity by one step (e.g., Moderate $→$ Minor).
  ℹtem \textbf{Boots of Sure Step (3 XP):} Ignore difficult terrain penalties.
\end{itemize}

⫕section{Removing an Enchantment}
An enchantment can be removed with a ritual (Wits + Arcana, DV 4) costing 1 XP per enchantment level. The item is not damaged.

\subsection{Weapon and Armor Acquisition}

\begin{itemize}
  ℹtem Mundane weapons and armor are acquired through play (found, looted, purchased with Supply, or granted as part of character background). No XP cost.
  ℹtem Masterwork or enchanted items require XP expenditure (as Assets) or narrative quests.
  ℹtem A character begins play with one Light or Medium weapon and one set of Light or Medium armor appropriate to their background.
\end{itemize}

